\section{Methods}\label{sec:methods}

In this section, we detail the construction of \repqa, the evaluation metrics, the readability-enhancement strategies we benchmarked, the study design behind our three research questions, and the experimental setup.

\subsection{Dataset construction}\label{sec:dataset_curation}

\paragraph{Expert-grounded patient education QA.}
We collected a diverse set of publicly available healthcare-related QA pairs from 27 authoritative online resources covering a broad range of public health topics, including general health, cardiovascular disease, cholesterol, diabetes, HIV, diet, exercise, alcohol, smoking, sleep, mental health, sexual health, and health equity.
Sources were selected for their accessibility to lay audiences and their coverage of frequently asked questions by the general public.
For the evaluation dataset, domain experts further manually filtered out QA pairs containing overly technical content, following national, authoritative research and clinical guidelines~\citep{doody2016nursing,chiang2016guidance,corry2013developing}.
Questions are phrased in plain language and do not include user profiles.

The resulting evaluation dataset includes 533 QA pairs spanning four categories, i.e., \textit{Suggestion}, \textit{Fact}, \textit{Definition}, and \textit{Rationale}, as summarized in Table~\ref{tab:dataset_category}:
\begin{itemize}[leftmargin=*]
\item \textit{Suggestion}: Questions seeking advice. These responses must be highly readable to ensure that users can easily follow recommendations without ambiguity.

\item \textit{Fact}: Questions requesting objective, verifiable information. Presenting facts in lay-accessible language is essential for avoiding misinterpretation.

\item \textit{Definition}: Questions asking for explanations of medical terms or conditions. Definitions require high readability, as they often introduce unfamiliar terminology.

\item \textit{Rationale}: Questions seeking justifications behind medical guidance. These responses must balance clarity with accuracy in plain language to build user understanding.
\end{itemize}

\paragraph{Large-scale healthcare questions.}
In addition to the evaluation QA pairs, we curated a large-scale collection consisting of 36,060 public health questions.
These queries span an extensive range of public health topics, including sleep, exercise, smoking, weight, diet, HIV care, mental health, cardiovascular disease, diabetes, cholesterol, and SOGI, drawn from 52 publicly available resources.
The data sources were manually selected by nursing and public health professionals, following the same guidelines used in curating the evaluation dataset.
Summary statistics and representative examples for each category are presented in Table~\ref{tab:training_dataset}. For each question, the reference answer was generated by GPT-4o-mini using supporting materials as context. During post-training, this context was not used to train the model’s parametric knowledge. The datasets were subsequently reviewed by domain experts to ensure quality.


\begin{table*}[t]
    \centering
    \caption{Examples and distribution of question categories in the collected training datasets. A single question can be categorized into multiple categories.}
    \resizebox{\textwidth}{!}{
    \begin{tabular}{clc}
    \toprule
    Category & Example & Proportion \\
    \midrule
    \midrule
    Sleep & How does sleep apnea during pregnancy potentially affect newborns? & 8.1\% \\
    Exercise & What are some effective strategies for improving family health through physical activities? & 6.3\% \\
    Smoking & What are some effective strategies individuals can use to quit smoking and support a smoke-free lifestyle? & 4.5\% \\
    Weight & How is body mass index (BMI) calculated, and what factors can affect its accuracy? & 7.2\% \\
    Diet & How can someone maintain a healthy diet while dining out? & 10.3\% \\
    HIV Care & How can gay men effectively manage their risk of HIV? & 14.4\% \\
    Mental Health & How can parents support the mental health of their children? & 12.6\% \\
    Cardiovascular & What is the primary difference between cardiac arrest and a heart attack? & 18.7\% \\
    Diabetes & How do BMI categories relate to the risk of type 2 diabetes across different age groups? & 5.2\% \\
    Cholesterol & What is the significance of LDL and HDL cholesterol in relation to heart health? & 5.1\% \\
    SOGI & How can schools and communities support LGBTQI+ youth to create a safer environment? & 12.4\% \\
    Others & What are the potential benefits of utilizing specialized healthcare models in patient treatment? & 33.7\% \\
    \bottomrule
    \end{tabular}
    }
    \label{tab:training_dataset}
\end{table*}

\subsection{Evaluation metrics}\label{subsec:metrics}

We evaluated responses with a suite of four complementary metrics: three readability metrics, i.e., the Flesch--Kincaid grade level (FK), the professional-language proportion (PL), and the Yngve depth (YD), and one informativeness metric, the informative cosine (IC) between a response and its reference answer. The three readability metrics capture word difficulty, domain-specific jargon, and syntactic complexity, respectively, while IC measures how much of the reference answer's content the response conveys.

\paragraph{Flesch-Kincaid grade level (FK).}
We used the Flesch-Kincaid grade level~\citep{kincaid1975derivation} as one of our readability metrics, an established standard for assessing reading difficulty that ranges from kindergarten to post-graduate levels.
The calculation is based on the average number of syllables in a word and the average number of words in a sentence:
\begin{equation}
\label{equ:grade}
    \pi = 0.39\left(\frac{\text{\#Words}}{\text{\#Sentences}}\right) + 11.8\left(\frac{\text{\#Syllables}}{\text{\#Words}}\right) - 15.59
\end{equation}
The intuition is that sentences with many words, or words with numerous syllables, are generally more difficult to read, corresponding to a higher Flesch-Kincaid grade level.
The range of grade levels can be categorized into different school levels, e.g., 6--9 corresponds to the middle school level.

\paragraph{Professional language (PL).}
During communication between the agent and lay users without a medical background, the frequency of professional terms can significantly affect the user's reading experience.
In light of this, we incorporated the proportion of health-related professional terms in the text as one of our readability metrics:
\begin{equation}
\label{equ:prof}
    \rho = 100\% \cdot \frac{\text{Total Professional Terms}}{\text{Total Words}}
\end{equation}
We compiled a list of professional terms from the Harvard Medical Dictionary\footnote{\url{https://www.health.harvard.edu/a-through-c}}, resulting in a vocabulary of 2,058 terms.
This metric complements the Flesch-Kincaid grade level by incorporating domain-specific terminology, which may affect comprehension for lay audiences but is not captured by traditional readability formulas. A lower PL indicates fewer jargon terms and is preferable for lay readers.

\paragraph{Yngve depth (YD).}
The Flesch--Kincaid grade and professional language capture word- and term-level difficulty, but not the syntactic complexity of a response, which also affects how easily lay users can follow it.
To capture this dimension, we adopted the Yngve depth~\citep{yngve1960model}, a measure of syntactic complexity derived from the constituency parse of a sentence.
Concretely, we parse each sentence and, for every word, accumulate the left-branching load along its path from the root, i.e., the number of as-yet-unexpanded sibling constituents to its right at each step; YD is the mean of these per-word depths, averaged over sentences.
Deeper, more nested, or more heavily left-branching structures yield a higher Yngve depth, so a lower YD corresponds to simpler, easier-to-follow syntax.

\paragraph{Informative cosine (IC).}
Readability metrics alone do not penalize a response that is easy to read but omits the essential content, so we additionally measured how much of the reference answer's information a response conveys.
Each evaluation question is paired with an expert-grounded reference answer.
We embedded both the model response and the reference answer with a sentence encoder (\texttt{all-MiniLM-L6-v2}~\citep{wang2020minilm}) and computed the cosine similarity between the two embeddings:
\begin{equation}
\label{equ:ic}
    \text{IC} = \frac{\mathbf{e}_{\text{resp}} \cdot \mathbf{e}_{\text{ref}}}{\lVert \mathbf{e}_{\text{resp}} \rVert \, \lVert \mathbf{e}_{\text{ref}} \rVert}
\end{equation}
A higher IC indicates that the response is more semantically aligned with the reference answer, and thus more informative. An ideal healthcare agent should produce responses that are both easy to understand (low FK toward the target, low PL, low YD) and informative (high IC).

\subsection{Readability-enhancement strategies}\label{sec:strategies}

Beyond evaluating off-the-shelf models, we benchmarked five strategies for steering responses toward the target reading level, grouped into training-free controls and reinforcement-learning fine-tuning. The first three are applied to a frozen base model and require no parameter updates:
\begin{itemize}[leftmargin=*]
\item \textbf{Prompt}: instruction-based control, in which the base model is prompted to answer in simple, plain language at roughly a 5th-grade level, using short sentences and minimal medical jargon~\citep{ribeiro-etal-2023-generating,malik2024tarzan}.
\item \textbf{CoT}: chain-of-thought prompting~\citep{wei2022chain}, which instructs the model to reason step by step before producing its final answer; only the final answer is scored.
\item \textbf{Lookahead}: readability-guided lookahead decoding~\citep{ribeiro-etal-2023-generating}. At each decoding step, the top candidate tokens are each extended by a short greedy rollout, the rollouts are scored by how close their Flesch--Kincaid grade is to the target, and the next-token log-probabilities are biased toward the candidate whose readability is closest to the target.
\end{itemize}
The remaining two fine-tune the base model with readability as a verifiable reward (Eq.~\ref{equ:reward}): \textbf{GRPO}~\citep{shao2024deepseekmath} and \textbf{DAPO}~\citep{yu2025dapo}, where DAPO improves training stability over GRPO via decoupled clipping and dynamic sampling. We applied these strategies to three open-source base models (LLaMA3.1-8B, Qwen3-8B, and Qwen3-4B).

\subsection{Study design}\label{sec:questions}

Using \repqa, we studied LLMs' ability to adapt, assess, and optimize readability. We posed three research questions.

\paragraph{RQ1: Can LLMs generate responses at a specified reading level?}
People with different educational backgrounds have varying health literacy.
Therefore, healthcare agents should adapt the readability of their answers to a requested level (e.g., 5th grade or middle school).
We examined whether LLMs follow explicit readability instructions without sacrificing informativeness.
We compared each response's Flesch–Kincaid grade to the requested level and assessed informativeness via the informative cosine (IC).
In addition, we evaluated LLMs' readability-understanding capability: LLMs were asked to predict the reading level of given responses, and we quantified their accuracy against the Flesch–Kincaid grade.

\paragraph{RQ2: Can LLMs determine what readability is needed for a given question?}
Oversimplifying medical content omits qualifiers and loses information, while overly technical language reduces accessibility. Thus, selecting an appropriate reading level for each question is critical.
We investigated whether an LLM can infer the target level from the question alone (without user metadata).
Concretely, we extended each question in \repqa into three variants that differ in required medical knowledge, and evaluated the accuracy of the models' predictions.

\paragraph{RQ3: Can readability serve as a reward signal?}
Reinforcement learning has shown promise for steering LLMs toward preferred behaviors. We studied whether readability can act as a useful reward to drive responses toward a requested level without sacrificing content quality.
Specifically, the reward combines the evaluation metrics, i.e., the three readability metrics and the informativeness metric, so that the policy is rewarded for being both easy to read and faithful to the reference answer:
\begin{equation}
\label{equ:reward}
r = -\Big(\alpha\,|\pi - 6| + \beta\,\rho + \gamma\,\max(0,\, \text{YD} - \delta_0) + \eta\,\max(0,\, L - L_{\max})\Big) + \mu\,(\text{IC} - 1),
\end{equation}
where $\pi$ is the Flesch--Kincaid grade (Eq.~\ref{equ:grade}), $\rho$ the professional-language proportion (Eq.~\ref{equ:prof}), $\text{YD}$ the Yngve depth, and $\text{IC}$ the informative cosine (Eq.~\ref{equ:ic}); $\alpha$, $\beta$, $\gamma$, $\eta$, and $\mu$ are weights that balance the components.
A Flesch--Kincaid grade closer to $6$ (the middle-school reading level), a lower professional-language proportion, a lower Yngve depth above the baseline $\delta_0$, and a higher informative cosine (closer to the reference) all increase the reward, while an auxiliary penalty discourages responses longer than a token budget $L_{\max}$.
The informativeness term subsumes the trivial-answer penalty used in prior formulations: a content-free response has low similarity to the reference and is therefore penalized automatically.
We optimized this reward with the two reinforcement-learning algorithms described in Section~\ref{sec:strategies} (GRPO and DAPO), and compared them against the training-free controls (Prompt, CoT, and Lookahead).

\subsection{Experimental setup}\label{sec:exp_setup}

We presented an evaluation of 19 LLMs, including 5 proprietary models and 14 open-source models.
For proprietary models, we considered the GPT-5 series (GPT-5.1, GPT-5.4, and GPT-5.5)~\citep{openai2025gpt5} and GPT-4.1 (including GPT-4.1-mini)~\citep{openai2024gpt4technicalreport,openai2024gpt4ocard}.
For open-source models, we used the LLaMA~3.1 series~\citep{grattafiori2024llama3herdmodels}, the Qwen-2.5 and Qwen-3 series~\citep{qwen2.5}, Phi-4~\citep{abdin2024phi4technicalreport}, Mistral~\citep{jiang2024mixtralexperts}, BioMistral~\citep{labrak2024biomistralcollectionopensourcepretrained}, InternLM3~\citep{cai2024internlm2}, Gemma-3~\citep{gemmateam2025gemma3}, and Yi-1.5~\citep{ai2025yiopenfoundationmodels}.
We used their instruction-tuned versions by default.